% Generated draft: requires mathematical and source review.
\documentclass[11pt]{article}
\usepackage[margin=25mm]{geometry}
\usepackage{fontspec}
\setmainfont{Latin Modern Roman}
\usepackage{amsmath,amssymb,mathtools,unicode-math}
\setmathfont{Latin Modern Math}
\usepackage{enumitem,needspace,xurl,hyperref}
\setlength{\parindent}{0pt}
\setlength{\parskip}{5pt}
\setlength{\emergencystretch}{3em}
\newcommand{\N}{\mathbb{N}}
\newcommand{\Z}{\mathbb{Z}}
\newcommand{\Q}{\mathbb{Q}}
\newcommand{\R}{\mathbb{R}}
\newcommand{\C}{\mathbb{C}}
\newcommand{\F}{\mathbb{F}}
\newcommand{\A}{\mathbb{A}}
\newcommand{\PP}{\mathbb{P}}
\newcommand{\E}{\mathbb{E}}
\newcommand{\eps}{\varepsilon}
\newcommand{\dd}{\,\mathrm{d}}
\begin{document}
\section*{UnsolvedMath ID: 30001372}
\textbf{Source problem code:} catalogue entry
\subsection{Definitions and mathematical statement}
The supplied question concerns arbitrary commutative rings, with no Noetherian hypothesis. For reference only, under the convention that a ring \(R\) has an identity, an element \(a\in R\) is regular if multiplication \(R\to R\), \(r\mapsto ar\), is injective. One standard convention calls a finite sequence \(a_1,\ldots,a_n\in R\), with \(n\geq 1\), an \(R\)-regular sequence if \((a_1,\ldots,a_n)\neq R\) and, for every \(i\in\{1,\ldots,n\}\), multiplication by the image of \(a_i\) on \(R/(a_1,\ldots,a_{i-1})\) is injective; here the ideal generated by the empty list is \((0)\). These conventions are explanatory: the supplied record does not establish that they are the conventions intended in the report.

The supplied question is: \emph{How can one avoid reliance on the Regular Element Theorem when developing a theory of regular sequences over arbitrary commutative rings without Noetherian hypotheses?}

The statement of the named Regular Element Theorem, the particular results that the requested theory must establish, and the permitted meaning of avoiding reliance on that theorem are absent. Consequently, the supplied material specifies a methodological question but does not determine a precise theorem, construction, or quantified existence assertion. In particular, the displayed terminology must not be treated as resolving the missing source context.
\subsection{Short English statement}
The question asks how to develop regular-sequence theory for arbitrary commutative rings without assuming Noetherianity or relying on the named Regular Element Theorem. The theorem and the required scope of the theory are unspecified, so a precise mathematical target requires further source context.
\subsection{Sources}
\begin{itemize}
\item Supplied local UnsolvedMath registry record: permanent ID \(30001372\), title \emph{Regular Sequences without the Regular Element Theorem}, exact source problem code \texttt{catalogue entry}, in the Oberwolfach Reports Open Problems collection. The record supplies proposed year 2009 and names no proposer.
\item Supplied source citation: \emph{Mini-Workshop: Formal Methods in Commutative Algebra: A View Toward Constructive Homological Algebra} (2010), Oberwolfach Reports, pp.\ 2731--2770, DOI \texttt{10.4171/owr/2009/50}, \url{https://doi.org/10.4171/owr/2009/50}. The background describes Oberwolfach Reports as documenting workshops and their open-problem discussions. No specific page or statement of the Regular Element Theorem is supplied.
\item Supplied literature source, labelled ``original Oberwolfach report 4133'': \url{https://doi.org/10.4171/owr/2009/50}. The literature assessment dated \texttt{2026-08-21} labels the problem open and describes the extracted sentence as a context-dependent question about avoiding a preceding algebraic result, rather than a fully specified theorem. It reports that no resolution could be verified from the cited passage or identified literature. This is supplied dated curation, not an independently verified literature assessment or source visit; the record requests independent verification of time-sensitive status and cited claims.
\end{itemize}
\end{document}
